\section{Unfolding and the problem of finding constant series}
\label{sec:unfolding}

\subsection{What a volume computation determines}

Let $\Gamma$ be the mapping class group of an oriented finite-type surface
with fixed boundary lengths, and let $\gamma$ be a fixed nonperipheral
curve or multicurve class. Use the same WP orbifold measure on both
quotients. For a measurable nonnegative function, or one for which
$\int_{\Gamma\backslash\T}\sum_\delta|F(\ell_\delta)|\,\dV<\infty$,
the basic unfolding identity is
\begin{equation}
 \int_{\Gamma\backslash\T}
 \sum_{\delta\in\Gamma\gamma}F(\ell_\delta(X))\,\dV(X)
 =\int_{\Stab(\gamma)\backslash\T}
             F(\ell_\gamma(X))\,\dV(X).
 \label{eq:unfold}
\end{equation}
The orbit contains each class once; $\Stab(\gamma)$ preserves the entire
chosen decoration, including any orientation, ordering or component
weights. Integration over coset translates of a fundamental domain proves
the formula, with ineffective isotropy retained on both sides. For
nonnegative functions either side may be infinite. Simplicity of the
orbit curves is not required.

Simplicity becomes important when evaluating the right side. If $\gamma$
is a disjoint simple multicurve, cut along it and use Fenchel--Nielsen
coordinates. Each twist period integrates to $\ell_j$, so the resulting
integral has factors $\prod_j\ell_j\,\dd\ell_j$ times WP volumes of the
components, with the relevant symmetry factors. This is the geometric
integration mechanism underlying Mirzakhani's volume recursion
\cite[Theorem 7.1 and Lemma 7.3]{MirzVolumes}. The symmetry and
one-handle factors in that formula must be converted using the volume
convention for each component; they are not extra factors to append to
an already stack-normalized quotient.

Suppose a series $S(X)$ is independently known to be constant, say $C$,
and $0<V_{g,n}<\infty$ in the chosen normalization.
Then
\begin{equation}
 C=\frac{1}{V_{g,n}}
       \int_{\M_{g,n}}S(X)\,\dV(X),
 \label{eq:recoverconstant}
\end{equation}
and unfolding can calculate its value. If constancy has not been proved,
the same quotient computes only its average. This is the precise logical
distinction behind the proposal to recover an identity from a volume.

On the punctured torus, for an unoriented essential simple curve orbit,
the stack-normalized formula is
\begin{equation}
 \int_{\M_{1,1}}\sum_{\alpha\text{ simple}}f(\ell_\alpha)\,\dV
       =\frac12\int_0^\infty x f(x)\,\dd x.
 \label{eq:torusunfold}
\end{equation}
Indeed, the effective stabilizer quotient has coordinates $x>0$ and
$0\leq\tau<x$, while the ineffective torus involution contributes the
single factor $1/2$. This is the stack convention of
Section~\ref{sec:torussphere}; it is not a second quotient by cuff exchange.
For $f(x)=(1+e^x)^{-1}$, expand in $e^{-x}$ to obtain
\[
 \int_0^\infty\frac{x\,\dd x}{1+e^x}
   =\sum_{m\ge1}\frac{(-1)^{m-1}}{m^2}=\frac{\pi^2}{12}.
\]
Termwise integration is justified by
$\sum_m\int_0^\infty xe^{-mx}\,\dd x=\sum_m m^{-2}<\infty$.
For a finite numerical check, the integral beyond $R$ is at most
$(R+1)e^{-R}$, and the alternating sum through $m=M$ has error at most
$(M+1)^{-2}$.
The integral of the series is $\pi^2/24$; division by
$V_{1,1}=\pi^2/12$ gives $C=1/2$. Its pointwise constancy is McShane's
theorem~\cite{McShane}:
\begin{equation}
 \sum_{\alpha\text{ simple, unoriented}}
                 \frac{1}{1+e^{\ell_\alpha(X)}}=\frac12.
 \label{eq:mcshane}
\end{equation}
The tetrahedral volume calculation and the covering argument provide an
independent route to the volume used in this normalization check.

\subsection{Operations that preserve constancy}

Finite linear combinations of constant series preserve constancy.
Parameter integrals do too when the constant is integrable and Fubini's
theorem applies to the joint sum and integral, or when all terms are
nonnegative and Tonelli's theorem applies. For differentiation on an
identified parameter family, one sufficient hypothesis is local uniform
absolute convergence of both the series and its first derivatives.
Uniform convergence of the original series alone is insufficient.
Products give constant double series when Tonelli's theorem or absolute
convergence applies. These operations construct kernels together with
their required convergence hypotheses.

Their geometric interpretation needs to be retained. A product usually
indexes pairs of curves that may intersect; it is not automatically a
disjoint-multicurve sum. A parameter derivative also differentiates lengths
if the metric varies with that parameter. In particular, differentiating
a constant series with respect to the metric gives the one-form identity
\begin{equation}
 \sum_\gamma F'(\ell_\gamma)\,\dd\ell_\gamma=0,
 \label{eq:derivativeidentity}
\end{equation}
under the stated local $C^1$ control, for a kernel with no additional
explicit metric dependence. It does not generally give
another scalar identity by differentiating only the displayed weight.

Two possible tests for a proposed constant series follow. One can try to
show that~\eqref{eq:derivativeidentity} vanishes on a spanning set of tangent
directions at every point; a $C^1$ series is then constant on each
connected component. Alternatively, for real-valued $S\in L^2$ on a
finite positive-volume moduli space, let $\overline S$ be its mean. Then
\begin{equation}
 \int_{\M}(S-\overline S)^2\,\dV=0
 \label{eq:variance}
\end{equation}
proves constancy almost everywhere, and everywhere if $S$ is continuous,
since WP measure has full support. For a complex-valued series use
$|S-\overline S|^2$ instead of the square.
The second-moment calculation involves pairs of curves, including their
intersection patterns. It can be substantially harder than the first
unfolding integral.

Thus a useful search for further identities should ask for more than
formal closure: a natural index set, a manageable positive weight, a
geometric partition or telescoping mechanism, and a convergence proof.
The constructions in Sections~\ref{sec:figures} and~\ref{sec:exacttwist}
illustrate these requirements.

\subsection{Other established constant series}

There are also constant series with substantially different kernels.
For example, write
\[
 \mathcal L(x)=\operatorname{Li}_2(x)+\tfrac12\log x\log(1-x)
\]
for the Rogers dilogarithm on $(0,1)$, extended continuously to its
endpoints, and define
\begin{align*}
 D(r)={}&\mathcal L\bigl(\sech^2(r/2)\bigr)\\
 &+2\mathcal L\left(\frac{1+e^{-r}}2\right)
  -2\mathcal L\left(\frac{1-e^{-r}}2\right).
\end{align*}
On a complete finite-area once-punctured torus, Hu--Tan's identity is
\begin{equation}
 \sum_{\alpha\text{ simple, unoriented}}D(\ell_\alpha)=\frac{\pi^2}{2}.
 \label{eq:dilogtorus}
\end{equation}
Here the index set is the nonperipheral simple spectrum with its curve
multiplicities~\cite[Theorem 2, equation (2)]{HuTan}. Their Theorem 1
treats a positive geodesic boundary with a different, boundary-dependent
kernel; the displayed $D$ is the cusped specialization.

For a closed oriented surface of genus $g\ge2$, Luo--Tan give explicit
dilogarithmic weights $f(P),g(T)$ on properly embedded geometric pairs of
pants and one-holed tori, each subsurface counted once,
such that
\begin{equation}
 \sum_P f(P)+\sum_T g(T)=8\pi^2(g-1)=\operatorname{Vol}(T^1X).
 \label{eq:luotan}
\end{equation}
Their construction partitions unit tangent directions up to a null set
\cite[Theorem 1.1 and Section 1.2]{LuoTan}. Angular measure $2\pi$ on
each unit tangent circle, together with area $4\pi(g-1)$, explains the
normalization on the right.
It illustrates another source of constant series: the invariant measure
being partitioned need not live only on a horocycle or a geodesic boundary.
The precise, longer subsurface weights are given in that reference.
